\section{Dual Problem}
\label{sec-dual}
\subsection{Discrete dual}
\label{sec-discrete-dual}
The Kantorovich problem~\eqref{eq-kanto-discr} is a linear program so that one can equivalently compute its value by solving a dual linear program.

\begin{defn}[Admissible potentials]\label{def-admissible-potentials}
For a discrete problem with marginal sizes $n,m$ and cost matrix $\C\in\RR^{n\times m}$, a pair $(\fD,\gD)\in\RR^n\times\RR^m$ is admissible if it lies below the cost:
\eql{\label{eq-feasible-potential}
\PotentialsD(\C) \eqdef \enscond{
(\fD,\gD) \in \RR^n \times \RR^m
}{ \foralls (i,j) \in \range{n} \times \range{m}, \fD_i+\gD_j \leq \C_{i,j} }.
}
Equivalently, $\fD\oplus\gD\leq\C$ entrywise. In particular, the feasible set depends on the cost matrix, not on the marginal weights.
\end{defn}
\begin{prop}[Discrete Kantorovich duality]\label{prop-duality-discr}
One has
\eql{\label{eq-dual}
\MKD_\C(\a,\b) =
\umax{(\fD,\gD) \in \PotentialsD(\C)} \dotp{\fD}{\a} + \dotp{\gD}{\b}.
}
\end{prop}
\begin{proof}
For completeness, introduce multipliers $(\fD,\gD)$ for the two marginal constraints. The Lagrangian is
\eql{\label{eq-mk-lagr}
\mathcal L(\P,\fD,\gD)
=\dotp{\C}{\P}+\dotp{\a-\P\ones_m}{\fD}
+\dotp{\b-\P^\top\ones_n}{\gD}.
}
For fixed potentials, the dual function is \(\inf_{\P\geq0}\mathcal L(\P,\fD,\gD) =\dotp{\a}{\fD}+\dotp{\b}{\gD} +\inf_{\P\geq0}\dotp{\C-(\fD\oplus\gD)}{\P}.\)
The last infimum is
\[
\inf_{\P\geq0}\dotp{\Q}{\P}
=
\begin{cases}
0 & \text{if }\Q\geq0,\\
-\infty & \text{otherwise},
\end{cases}
\]
and is therefore finite exactly when $(\fD,\gD)\in\PotentialsD(\C)$. The primal polytope is nonempty because it contains the product coupling $\a\otimes\b$, and it is compact; hence the primal value is finite and attained. Finite-dimensional linear-programming strong duality then identifies this value with the maximum of the dual function and also gives dual attainment.
\end{proof}
\begin{prop}[Discrete complementary slackness]\label{prop-discrete-complementary-slackness}
For every $\P\in\CouplingsD(\a,\b)$ and $(\fD,\gD)\in\PotentialsD(\C)$, \(\dotp{\C}{\P}-\dotp{\fD}{\a}-\dotp{\gD}{\b} = \sum_{i,j}\P_{i,j}\bigl(\C_{i,j}-\fD_i-\gD_j\bigr) \geq0.\)
The plan and potentials are respectively primal and dual optimal if and only if
\eql{\label{eq-mk-pd-rel}
\Supp(\P)
\subset
\enscond{(i,j)\in\range{n}\times\range{m}}{\fD_i+\gD_j=\C_{i,j}}.
}
\end{prop}
\begin{proof}
The marginal constraints give
$\dotp{\fD}{\a}+\dotp{\gD}{\b}
=\sum_{i,j}\P_{i,j}(\fD_i+\gD_j)$, which proves the identity and nonnegativity. If both feasible objects are optimal, their values agree by Proposition~\ref{prop-duality-discr}; every nonnegative summand with $\P_{i,j}>0$ must therefore vanish. Conversely, the support condition makes the gap zero, so weak duality forces both objects to be optimal.
\end{proof}
\subsection{Auction Algorithm and Dual Prices}
\label{sec-auction-dual-ascent}
\paragraph{Bidding dynamics.}
Consider a square assignment problem of size $n\geq2$ with costs $C_{i,j}$ and rewrite it as the profit maximization problem with $a_{i,j}=-C_{i,j}$. The case $n=1$ is immediate.

\(v_i=\max_j (a_{i,j}-p_j),\qquad j_i\in\operatorname*{argmax}_j(a_{i,j}-p_j),\qquad w_i=\max_{j\neq j_i}(a_{i,j}-p_j).\)
\[
p_{j_i}\leftarrow p_{j_i}+v_i-w_i+\epsilon.
\]
\paragraph{Approximate optimality.}
For fixed prices $p$, eliminating the bidder utilities $u_i$ in the dual minimization gives the convex objective

\(D(p)=\sum_j p_j+\sum_i \max_j(a_{i,j}-p_j),\)
\begin{defn}[$\epsilon$-complementary slackness]\label{def-auction-eps-cs}
An assignment $\sigma$ and prices $p$ satisfy $\epsilon$-complementary slackness if, for every source $i$, \(a_{i,\sigma(i)}-p_{\sigma(i)} \geq \max_j(a_{i,j}-p_j)-\epsilon.\)
For a partial assignment, the same condition is required only for currently assigned sources.
\end{defn}
\begin{prop}[Auction invariant and finite termination]\label{prop-auction-termination}
At every iteration of Algorithm~\ref{alg:auction-bidding}, each currently owned pair satisfies $\epsilon$-complementary slackness. Moreover, if
\(\Delta_A\eqdef\max_{i,j}a_{i,j}-\min_{i,j}a_{i,j},\)
then the algorithm terminates after at most
\[
1+\left\lfloor\frac{n(\Delta_A+\epsilon)}{\epsilon}\right\rfloor
\]
bids.
\end{prop}
\begin{proof}
Immediately after source $i$ bids for $j_i$, its reduced profit there is \(a_{i,j_i}-p_{j_i}^{\mathrm{new}}=w_i-\epsilon,\)
whereas every alternative has reduced profit at most $w_i$. Later bids only increase competing prices. The price of $j_i$ itself cannot change while $i$ remains its owner, because such a change would reassign $j_i$ and make $i$ unassigned. Thus every owned pair satisfies the claimed invariant.

Each bid raises one price by at least $\epsilon$. After any nonterminal bid, at least one target different from the selected target remains unassigned; it has never been owned and still has price zero. Denote such a target by $k$. Since it was also available before the bid, the second-best value satisfies $w_i\geq a_{i,k}$, and hence
\(p_{j_i}^{\mathrm{new}}=a_{i,j_i}-w_i+\epsilon \leq \Delta_A+\epsilon.\)
Since this sum grows by at least $\epsilon$ per bid, there can be at most $n(\Delta_A+\epsilon)/\epsilon$ nonterminal bids, followed by at most one terminal bid.
\end{proof}
\begin{prop}[Auction optimality certificate]\label{prop-auction-eps-cs}
If a complete assignment $\sigma$ satisfies $\epsilon$-complementary slackness, then it is $n\epsilon$-optimal for the profit maximization problem, or equivalently $n\epsilon$-optimal for the original cost minimization problem. If all costs are integers and $\epsilon<1/n$, then $\sigma$ is optimal.
\end{prop}
\begin{proof}
Let $\tau$ be any assignment. By $\epsilon$-complementary slackness, \(a_{i,\tau(i)}-p_{\tau(i)} \leq \max_j(a_{i,j}-p_j) \leq a_{i,\sigma(i)}-p_{\sigma(i)}+\epsilon.\)
Summing over $i$ cancels prices, because both $\sigma$ and $\tau$ are permutations: \(\sum_i a_{i,\tau(i)} \leq \sum_i a_{i,\sigma(i)}+n\epsilon.\)
Thus no assignment has profit more than $n\epsilon$ above that of $\sigma$. Since $a=-C$, the same statement says that the cost of $\sigma$ is at most $n\epsilon$ above the minimum cost. If the costs are integers, all assignment costs are integers; a gap strictly smaller than one therefore forces the gap to be zero.
\end{proof}
\subsection{General formulation}
\label{sec-dual-general}
\paragraph{Continuous duality.}
\begin{prop}[Kantorovich duality]\label{prop-kantorovich-duality-general}
Assume that $\Xx$ and $\Yy$ are compact metric spaces and that $c\in\Cc(\Xx\times\Yy)$. Then
\eql{\label{eq-dual-generic}
\MK_\c(\al,\be) =
\umax{(\f,\g) \in \Potentials(\c)}
\int_\X \f(x) \d\al(x) + \int_\Y \g(y) \d\be(y),
}
where the set of admissible dual potentials is
\eql{\label{eq-dfn-pot-dual}
\Potentials(\c) \eqdef \enscond{
(\f,\g) \in \Cc(\X) \times \Cc(\Y)
}{
\forall (x,y), \f(x)+\g(y) \leq \c(x,y)
}.
}
Here, $(\f,\g)$ is a pair of continuous functions, often called ``Kantorovich potentials''.
The same formula extends under the usual lower-semicontinuity and integrability assumptions, replacing maxima by suprema when dual optimizers need not exist.
\end{prop}
\begin{proof}
Weak duality is immediate: if $f(x)+g(y)\leq c(x,y)$ and $\pi\in\Couplings(\al,\be)$, then
\(\int f\d\al+\int g\d\be = \int (f(x)+g(y))\d\pi(x,y) \leq \int c\d\pi.\)
Taking the supremum over admissible potentials and the infimum over couplings gives the first inequality.

For the reverse inequality, write $V=\MK_\c(\al,\be)$ and consider the convex cone
\[
\mathcal K
\eqdef
\left\{\left(\pi_1,\pi_2,\int c\d\pi+r\right):
\pi\in\Mm_+(\X\times\Y),\ r\geq0\right\}
\]
in $\Mm(\X)\times\Mm(\Y)\times\RR$, equipped with the weak topologies on the measure factors. This cone is closed. Indeed, convergence of the first marginal implies a uniform bound on the masses of the underlying nonnegative measures. Compactness of $\X\times\Y$ then gives a convergent subnet $\pi_k\rightharpoonup\pi$; continuity of marginalization and of $\pi\mapsto\int c\d\pi$ identifies the limit with an element of $\mathcal K$, while the limiting epigraph remainder remains nonnegative.

For every $t<V$, the point $z_t=(\al,\be,t)$ does not belong to $\mathcal K$. Strong separation of this point from the closed convex cone gives a continuous linear functional. Because $\mathcal K$ is a cone containing the origin, the functional can be oriented so that, for some $F\in\Cc(\X)$, $G\in\Cc(\Y)$ and $\lambda\in\RR$,
\(\int F\d\al+\int G\d\be+\lambda t<0 \leq \int F\d\eta+\int G\d\zeta+\lambda s \qquad \text{for all }(\eta,\zeta,s)\in\mathcal K.\)
Testing $(0,0,r)\in\mathcal K$ gives $\lambda\geq0$, while testing the element generated by $\pi=\de_{(x,y)}$ and $r=0$ gives
$F(x)+G(y)+\lambda c(x,y)\geq0$. The multiplier cannot vanish: if $\lambda=0$, integrating this last inequality against $\al\otimes\be$ would contradict the strict inequality at $z_t$. Thus $\lambda>0$, and
\(f\eqdef-F/\lambda,\qquad g\eqdef-G/\lambda\)
satisfy $f\oplus g\leq c$ and
$\int f\d\al+\int g\d\be>t$. Letting $t\uparrow V$ proves the reverse inequality.

Take a maximizing sequence $(f_k,g_k)$. Successively replace it by \(\widetilde g_k(y)=\min_x\{c(x,y)-f_k(x)\}, \qquad \widetilde f_k(x)=\min_y\{c(x,y)-\widetilde g_k(y)\}.\)
Each replacement preserves feasibility and cannot decrease the objective. Fixing $x_0\in\X$, use the gauge freedom
$(\widetilde f_k,\widetilde g_k)\mapsto
(\widetilde f_k-a_k,\widetilde g_k+a_k)$
to impose $\widetilde f_k(x_0)=0$. Uniform continuity of $c$ shows that $\widetilde g_k$ inherits a common modulus in the second variable and $\widetilde f_k$ a common modulus in the first. The gauge gives a uniform bound on $\widetilde f_k$, and the first envelope then gives one on $\widetilde g_k$. Arzel\`a--Ascoli yields uniformly convergent subsequences. Their limit remains admissible and attains the dual supremum.
\end{proof}
\paragraph{Complementary slackness.}
\begin{prop}[Continuous complementary slackness]\label{prop-continuous-complementary-slackness}
For $\pi\in\Couplings(\al,\be)$ and $(f,g)\in\Potentials(\c)$, \(\int c\d\pi-\int f\d\al-\int g\d\be = \int\bigl(c(x,y)-f(x)-g(y)\bigr)\d\pi(x,y) \geq0.\)
The coupling and potentials are respectively primal and dual optimal if and only if this gap vanishes. In that case,
\eql{\label{eq-mk-pd-rel-cont}
\Supp(\pi)
\subset
\enscond{(x,y)\in\X\times\Y}{f(x)+g(y)=c(x,y)}.
}
\end{prop}
\begin{proof}
The marginal identities give the equality, and admissibility gives nonnegativity. A zero gap is equivalent to equality of the feasible primal and dual values, hence to joint optimality by Proposition~\ref{prop-kantorovich-duality-general}. Finally, the slack is continuous and nonnegative. If it were positive at a point of $\Supp(\pi)$, it would be bounded below by a positive constant on a neighborhood of positive $\pi$-mass, contradicting its zero integral.
\end{proof}
\(s(x,y) \eqdef c(x,y)-f(x)-g(y).\)
Dual feasibility is the statement $s\geq0$, while complementary slackness says that an optimal coupling can only live where this slack vanishes.

\subsection{\texorpdfstring{$c$-transforms}{c-transforms}}
\label{sec-c-transfo}
\paragraph{\texorpdfstring{Best-response potentials and the $c$-transform.}{Best-response potentials and the c-transform.}}
\(\usup{\g \in \Cc(\Yy)} \enscond{ \int g \d \be }{ \foralls (x,y), \g(y) \leq c(x,y) - \f(x) }.\)
\(\foralls y \in \Yy, \quad \g(y) \leq \f^\c(y)\)
\begin{defn}[$c$-transform]\label{def-c-transform}
For a real-valued function $f:\Xx\to\RR$, its $c$-transform is the extended-real-valued function
\eql{\label{eq-c-transform}
\foralls y \in \Y, \quad
\f^\c(y) \eqdef \uinf{x \in \X} \c(x,y) - \f(x).
}
For a real-valued function $g:\Yy\to\RR$, the $\bar c$-transform associated with $\bar c(y,x)=c(x,y)$ is \(\foralls x \in \X, \quad \g^{\bar\c}(x) \eqdef \uinf{y \in \Y} \c(x,y) - \g(y).\)
When $\Xx,\Yy$ are compact, $c$ is continuous and the input potential is continuous, these infima are finite minima and the transformed potential is continuous.
\end{defn}
Since $\be$ is nonnegative, maximizing $\int \g \d \be$ saturates the bound $\g\leq\f^\c$ on $\Supp(\be)$, equivalently $\be$-almost everywhere.

\begin{prop}[$c$-transforms solve the semi-relaxed problems]\label{prop-c-transform-semi-relaxed}
Assume the compactness and continuity hypotheses of Proposition~\ref{prop-kantorovich-duality-general}. For fixed $f\in\Cc(\X)$, the maximizers of the dual objective over all $g\in\Cc(\Y)$ such that $f\oplus g\leq c$ are exactly the functions satisfying $g=f^c$ $\be$-almost everywhere. Equivalently, $f^c$ gives the value of the one-marginal primal problem
\[
\inf_{\substack{\pi\in\Mm_+^1(\X\times\Y)\\ \pi_2=\be}}
\int c(x,y)\d\pi(x,y)-\int f(x)\d\pi_1(x)
=
\int f^c(y)\d\be(y).
\]
Symmetrically, for fixed $g$, the maximizers over $f$ are the functions satisfying $f=g^{\bar c}$ $\al$-almost everywhere.
\end{prop}
\begin{proof}
The constraint $f(x)+g(y)\leq c(x,y)$ for all $x$ is equivalent, for each fixed $y$, to
\(g(y)\leq \inf_x c(x,y)-f(x)=f^c(y).\)
Since $\be$ is nonnegative, the largest possible value of $\int g\d\be$ is obtained by saturating this pointwise upper bound on the support of $\be$. The proof for $f=g^{\bar c}$ is identical after exchanging the two marginals.

For the primal formula, disintegrate any feasible $\pi$ as $\pi(\d x,\d y)=\pi_y(\d x)\be(\d y)$. Then
\[
\int c\d\pi-\int f\d\pi_1
=
\int\left(\int (c(x,y)-f(x))\d\pi_y(x)\right)\d\be(y)
\geq
\int f^c(y)\d\be(y).
\]
The argmin of $x\mapsto c(x,y)-f(x)$ is nonempty and compact. The measurable maximum theorem provides a Borel selector $x_\star(y)$ from these argmins. The coupling
$\pi(\d x,\d y)=\de_{x_\star(y)}(\d x)\be(\d y)$
then attains equality.
\end{proof}
\((f,g)\longmapsto(f,f^c),\)
\((f,f^c)\longmapsto(f^{c\bar c},f^c).\)
\begin{prop}[Modulus stability of $c$-transforms]\label{prop-c-transform-lipschitz}
Suppose that $\omega_\Y$ is a modulus such that
\(|c(x,y)-c(x,y')| \leq\omega_\Y\bigl(d_\Y(y,y')\bigr) \qquad\text{for all }x,y,y'.\)
Then every finite $f^c$ has the same modulus. Symmetrically, a uniform modulus $\omega_\X$ of $c$ in its first variable is inherited by $g^{\bar c}$. In particular, a uniform $L$-Lipschitz bound on the cost yields an $L$-Lipschitz bound on the corresponding transform.
\end{prop}
\begin{proof}
For each $x$, set $F_x(y)=c(x,y)-f(x)$ and $F(y)=f^c(y)=\inf_x F_x(y)$. Since all the functions $F_x$ share the modulus $\omega_\Y$,
\[
|F(y)-F(y')|
=
\left|\inf_x F_x(y)-\inf_x F_x(y')\right|
\leq
\sup_x |F_x(y)-F_x(y')|
\leq
\omega_\Y\bigl(d_\Y(y,y')\bigr).
\]
The proof in the first variable is identical.
\end{proof}
\paragraph{Euclidean case.}
The Euclidean quadratic cost is the model case where $c$-transforms become ordinary convex conjugates after removing the quadratic terms. This is the algebraic bridge between Kantorovich duality and Brenier maps.

\(\int q(x,y)\d\pi(x,y) = \frac12\int\norm{x}^2\d\al(x) +\frac12\int\norm{y}^2\d\be(y) -\int \dotp{x}{y}\d\pi(x,y).\)
\(f^c(y) = \inf_x -\dotp{x}{y}-f(x) = -(-f)^*(y), \qquad h^*(y)\eqdef\sup_x \dotp{x}{y}-h(x).\)
On the full Euclidean space, $f^{c\bar c}=-(-f)^{**}$; hence $c$-closed functions are negatives of lower semicontinuous convex functions, and the closure is the smallest upper semicontinuous concave majorant of $f$. On restricted compact domains, the same statement holds with the additional restriction that the supporting slopes belong to the opposite domain.

\paragraph{Why hard alternating optimization stops.}
\begin{prop}[Algebra of $c$-transforms]\label{prop-c-transform-algebra}
Writing $\f^{c\bar c}\eqdef(\f^c)^{\bar c}$, the following pointwise identities hold whenever the transforms are well defined:
\[
\textnormal{(i) } f \leq f' \Rightarrow f^{c}\geq f'^{c},
\qquad
\textnormal{(ii) } f^{\c\bar{\c}} \geq f,
\qquad
\textnormal{(iii) } g^{\bar{\c}\c} \geq g,
\qquad
\textnormal{(iv) } \f^{\c\bar{\c}\c}=\f^{\c}.
\]
\end{prop}
\begin{proof}
The order reversal in (i) follows directly from the minus sign in the definition. For every $x$ and $y$, \(f^c(y)=\inf_{x'}\{c(x',y)-f(x')\} \leq c(x,y)-f(x),\)
so $c(x,y)-f^c(y)\geq f(x)$. Taking the infimum over $y$ proves (ii), and exchanging the two spaces proves (iii). Finally, (ii) and the order reversal imply
$f^c\geq f^{c\bar c c}$, whereas (iii) applied to $g=f^c$ gives
$f^{c\bar c c}\geq f^c$. This proves (iv).
\end{proof}
\eql{\label{eq-iter-c-trans}
(f,g) \mapsto (f,f^c) \mapsto (f^{c\bar c},f^c) \mapsto (f^{c\bar c},f^{c\bar c c}) = (f^{c\bar c},f^c) \ldots
}
For the bilinear cost $c(x,y)=-xy$ on a compact interval, $f^{c\bar c}$ is the smallest concave majorant representable with slopes in the opposite interval. A hard transform removes non-concave oscillations in one step rather than producing a gradual ascent.